\setlength{\tabcolsep}{4pt}
\begin{tabular}{@{}p{0.15\textwidth}p{0.81\textwidth}@{}}
\toprule
\textbf{System (v1, v3)} & Vous êtes juriste, spécialiste du droit québécois des assurances et de la rédaction législative. Vous évaluez si une reformulation simplifiée d'une clause préserve exactement sa portée juridique. \\
\midrule
\textbf{User, v1 (headline)} & \parbox[t]{0.80\textwidth}{\raggedright Voici une clause originale et sa version simplifiée.\\ \phantom{x}\\ CLAUSE ORIGINALE:\\ \emph{[original clause]}\\ \phantom{x}\\ VERSION SIMPLIFIÉE:\\ \emph{[simplified version]}\\ \phantom{x}\\ Évaluez UNIQUEMENT la préservation du sens juridique: la version simplifiée dit-elle exactement ce que dit l'originale sur le plan du droit? Ignorez complètement la lisibilité, le style et la simplicité -- ils sont évalués séparément.\\ \phantom{x}\\ Barème (échelle ordinale de 1 à 10):\\ - Partez de 10 (portée juridique identique).\\ - Retirez un point par erreur juridique identifiée, parmi:\\   * OMISSION: un élément juridiquement opérant de l'originale a disparu (condition, exclusion, délai, montant, partie liée, exception).\\   * HALLUCINATION: la version simplifiée ajoute un contenu juridique absent de l'originale.\\   * INCOHÉRENCE: la version simplifiée contredit l'originale (modalité, polarité, partie, quantum, portée).\\   * CONFUSION: la version simplifiée rend la règle juridiquement ambiguë ou inintelligible alors que l'originale ne l'était pas.\\ - 1 = le sens juridique n'est pas préservé du tout (règle inversée, autre règle,\\   ou contenu sans rapport).\\ \phantom{x}\\ Répondez uniquement par un objet JSON, sans texte autour:\\ {"erreurs": [{"type": "...", "detail": "..."}], "score": <entier de 1 à 10>}} \\
\midrule
\textbf{User, v3 (bare)} & \parbox[t]{0.80\textwidth}{\raggedright Voici une clause originale et sa version simplifiée.\\ \phantom{x}\\ CLAUSE ORIGINALE:\\ \emph{[original clause]}\\ \phantom{x}\\ VERSION SIMPLIFIÉE:\\ \emph{[simplified version]}\\ \phantom{x}\\ Sur une échelle de 1 à 10, dans quelle mesure la version simplifiée préserve-t-elle le sens juridique de la clause originale? 10 = portée juridique identique, 1 = sens juridique pas préservé du tout. Ignorez la lisibilité et le style.\\ \phantom{x}\\ Répondez uniquement par un objet JSON, sans texte autour:\\ {"erreurs": [], "score": <entier de 1 à 10>}} \\
\midrule
\textbf{v2} & The v1 rubric translated into English, over the same French clauses (system prompt in English). Full text in the repository, \texttt{src/llm\_judge.py}. \\
\bottomrule
\end{tabular}
